% GENERATED FILE. Edit canonical lessons, then run npm run generate:books.
% canonical-source-hash: fnv1a64:ddae6838f3a07042
% canonical-lessons: ZH-C247-jingli, ZH-C247-yukuai, ZH-C247-techan, ZH-C247-jinianpin, ZH-C247-huiyi, ZH-R247-review-242-247, ZH-R247-review-final-247

\chapter{Back From a Trip}
\label{ch:zh-back-from-a-trip}
\hlchaptermodality{247}

\begin{chapteropening}
\textbf{By the end of this chapter:} \emph{I can talk about a trip: what I went through, how pleasant it was, local specialities, souvenirs and memories.}

Complete the last lesson of chapter 247: I can talk about a trip: what I went through, how pleasant it was, local specialities, souvenirs and memories.
\end{chapteropening}

\section[jīnglì]{jīnglì — an experience (something lived through)}

\label{lesson:ZH-C247-jingli}



\begin{quote}
Before the new one: say the Mandarin for a pond, then the Mandarin for to row a boat.
\end{quote}

\subsection*{You'll want to know: jīnglì}
\emph{jīnglì} — \textquotedblleft{}an experience (something lived through)\textquotedblright{}.

\begin{quote}
\emph{zhè cì lǚyóu shì yí cì hěn hǎo de jīnglì} — This trip was a great experience.
\end{quote}

Its characters come later; for now it is a word for the ear.

\begin{cousinweb}
\emph{jīnglì} is built from \emph{jīng} 'to pass through' and \emph{lì} 'to go through', both inherited from Middle Chinese, \emph{keng} and \emph{lek} (Guangyun, in Baxter's transcription), which glosses the second 'to pass' and 'to go'.
\end{cousinweb}

\begin{grammarlens}[title={What you've built}]
One more word for talking about the past.
\end{grammarlens}

\subsection*{Your turn}
\begin{itemize}
\raggedright
  \item \emph{Say it:} \emph{jīnglì}
  \item \emph{Say it:} \emph{jīnglì}, once more
  \item \emph{Say it:} \emph{zhè cì lǚyóu shì yí cì hěn hǎo de jīnglì} — This trip was a great experience.
  \item \emph{Recall:} say the Mandarin for a pond, then the Mandarin for to row a boat, then say \emph{jīnglì} again
\end{itemize}

\begin{tcolorbox}[breakable,colback=teal!4,colframe=teal!35!black,title={Before you move on}]
What does jīnglì mean? (\textquotedblleft{}An experience (something lived through)\textquotedblright{}.) Say it once more.
\end{tcolorbox}
\section[yúkuài]{yúkuài — pleasant; cheerful}

\label{lesson:ZH-C247-yukuai}



\begin{quote}
Before the new one: say the Mandarin for to row a boat, then the Mandarin for something lived through.
\end{quote}

\subsection*{You'll want to know: yúkuài}
\emph{yúkuài} — \textquotedblleft{}pleasant; cheerful\textquotedblright{}.

\begin{quote}
\emph{zhōumò wǒmen wán de hěn yúkuài} — We had a lovely time at the weekend.
\end{quote}

Its characters come later; for now it is a word for the ear.

\begin{cousinweb}
\emph{yúkuài} is built from \emph{yú} 'glad' and \zh{快} 'pleased', both inherited from Middle Chinese, \emph{yu} and \emph{khwaejH} (Guangyun, in Baxter's transcription), which glosses the first 'pleased' and 'happy' and the second 'to one's liking' and 'glad'.
\end{cousinweb}

\begin{grammarlens}[title={What you've built}]
One more word for talking about the past.
\end{grammarlens}

\subsection*{Your turn}
\begin{itemize}
\raggedright
  \item \emph{Say it:} \emph{yúkuài}
  \item \emph{Say it:} \emph{yúkuài}, once more
  \item \emph{Say it:} \emph{zhōumò wǒmen wán de hěn yúkuài} — We had a lovely time at the weekend.
  \item \emph{Recall:} say the Mandarin for to row a boat, then the Mandarin for something lived through, then say \emph{yúkuài} again
\end{itemize}

\begin{tcolorbox}[breakable,colback=teal!4,colframe=teal!35!black,title={Before you move on}]
What does yúkuài mean? (\textquotedblleft{}Pleasant; cheerful\textquotedblright{}.) Say it once more.
\end{tcolorbox}
\section[tèchǎn]{tèchǎn — a local speciality; a regional product}

\label{lesson:ZH-C247-techan}



\begin{quote}
Before the new one: say the Mandarin for something lived through, then the Mandarin for pleasant and cheerful.
\end{quote}

\subsection*{You'll want to know: tèchǎn}
\emph{tèchǎn} — \textquotedblleft{}a local speciality; a regional product\textquotedblright{}.

\begin{quote}
\emph{zhèr de tèchǎn hěn hǎochī} — The local specialities here are delicious.
\end{quote}

Its characters come later; for now it is a word for the ear.

\begin{cousinweb}
\emph{tèchǎn} is built from \emph{tè} 'special' and \emph{chǎn} 'to produce', both inherited from Middle Chinese, \emph{dok} and \emph{sreanX} (Guangyun, in Baxter's transcription), which glosses the first 'a bull' and 'alone' and the second 'to give birth'; the ch of chǎn is not the regular outcome of that form.
\end{cousinweb}

\begin{grammarlens}[title={What you've built}]
One more word for talking about the past.
\end{grammarlens}

\subsection*{Your turn}
\begin{itemize}
\raggedright
  \item \emph{Say it:} \emph{tèchǎn}
  \item \emph{Say it:} \emph{tèchǎn}, once more
  \item \emph{Say it:} \emph{zhèr de tèchǎn hěn hǎochī} — The local specialities here are delicious.
  \item \emph{Recall:} say the Mandarin for something lived through, then the Mandarin for pleasant and cheerful, then say \emph{tèchǎn} again
\end{itemize}

\begin{tcolorbox}[breakable,colback=teal!4,colframe=teal!35!black,title={Before you move on}]
What does tèchǎn mean? (\textquotedblleft{}A local speciality; a regional product\textquotedblright{}.) Say it once more.
\end{tcolorbox}
\section[jìniànpǐn]{jìniànpǐn — souvenir}

\label{lesson:ZH-C247-jinianpin}



\begin{quote}
Before the new one: say the Mandarin for pleasant and cheerful, then the Mandarin for a local speciality.
\end{quote}

\subsection*{You'll want to know: jìniànpǐn}
\emph{jìniànpǐn} — \textquotedblleft{}souvenir\textquotedblright{}.

\begin{quote}
\emph{wǒ mǎile liǎng ge jìniànpǐn} — I bought two souvenirs.
\end{quote}

Its characters come later; for now it is a word for the ear.

\begin{cousinweb}
\emph{jìniànpǐn} is built from \emph{jì} 'to record', \emph{niàn} 'to think of' and \emph{pǐn} 'an article', all inherited from Middle Chinese, \emph{kiX}, \emph{nemH} and \emph{phimX} (Guangyun, in Baxter's transcription), which glosses the second 'to think' and the third 'a kind' and 'a rank'.
\end{cousinweb}

\begin{grammarlens}[title={What you've built}]
One more word for talking about the past.
\end{grammarlens}

\subsection*{Your turn}
\begin{itemize}
\raggedright
  \item \emph{Say it:} \emph{jìniànpǐn}
  \item \emph{Say it:} \emph{jìniànpǐn}, once more
  \item \emph{Say it:} \emph{wǒ mǎile liǎng ge jìniànpǐn} — I bought two souvenirs.
  \item \emph{Recall:} say the Mandarin for pleasant and cheerful, then the Mandarin for a local speciality, then say \emph{jìniànpǐn} again
\end{itemize}

\begin{tcolorbox}[breakable,colback=teal!4,colframe=teal!35!black,title={Before you move on}]
What does jìniànpǐn mean? (\textquotedblleft{}Souvenir\textquotedblright{}.) Say it once more.
\end{tcolorbox}
\section[huíyì]{huíyì — to recall; to think back on}

\label{lesson:ZH-C247-huiyi}



\begin{quote}
Before the new one: say the Mandarin for a local speciality, then the Mandarin for a souvenir.
\end{quote}

\subsection*{You'll want to know: huíyì}
\emph{huíyì} — \textquotedblleft{}to recall; to think back on\textquotedblright{}.

\begin{quote}
\emph{wǒ chángcháng huíyì nà cì lǚyóu} — I often think back to that trip.
\end{quote}

Its characters come later; for now it is a word for the ear.

\begin{cousinweb}
\emph{huíyì} is built from \zh{回} 'to return' and \emph{yì} 'to remember'. The first is inherited from Middle Chinese \emph{hwoj} (Guangyun, in Baxter's transcription); the dictionary glosses the second, \emph{'ik} there, 'to think of'.
\end{cousinweb}

\begin{grammarlens}[title={What you've built}]
One more word for talking about the past.
\end{grammarlens}

\subsection*{Your turn}
\begin{itemize}
\raggedright
  \item \emph{Say it:} \emph{huíyì}
  \item \emph{Say it:} \emph{huíyì}, once more
  \item \emph{Say it:} \emph{wǒ chángcháng huíyì nà cì lǚyóu} — I often think back to that trip.
  \item \emph{Recall:} say the Mandarin for a local speciality, then the Mandarin for a souvenir, then say \emph{huíyì} again
\end{itemize}

\begin{tcolorbox}[breakable,colback=teal!4,colframe=teal!35!black,title={Before you move on}]
What does huíyì mean? (\textquotedblleft{}To recall; to think back on\textquotedblright{}.) Say it once more.
\end{tcolorbox}
\section[(dialogue)]{Looking back: from Farm Animals and Bugs to Back From a Trip}

\label{lesson:ZH-R247-review-242-247}



\begin{quote}
Say the Mandarin for a pig and a goat.
\end{quote}

\subsection*{Your turn}
Say these aloud:

\begin{itemize}
\raggedright
  \item a pig, a goat, a snake, a bug, a mosquito
  \item climate, a snowflake, to change over time, to freeze over, temperature
  \item a typhoon, a rain shower, who knows, in case, to carry on
  \item body temperature, to be catching, a face mask, mood, to look after yourself
  \item a village, a bridge, a stone, a pond, to row a boat
  \item something lived through, pleasant and cheerful, a local speciality, a souvenir, to think back on
\end{itemize}

\begin{tcolorbox}[breakable,colback=teal!4,colframe=teal!35!black,title={Before you move on}]
A pig? (\emph{zhū}.) A goat? (\emph{shānyáng}.) A snake? (\emph{shé}.) A bug? (\emph{chóngzi}.) A mosquito? (\emph{wénzi}.) Climate? (\textbf{\zh{气候}}, \emph{qìhòu}.) A snowflake? (\textbf{\zh{雪花}}, \emph{xuěhuā}.) To change over time? (\emph{biànhuà}.) To freeze over? (\emph{jiébīng}.) Temperature? (\emph{wēndù}.) A typhoon? (\emph{táifēng}.) A rain shower? (\emph{zhènyǔ}.) Who knows? (\emph{shuōbudìng}.) In case? (\emph{wànyī}.) To carry on? (\emph{jìxù}.) Body temperature? (\emph{tǐwēn}.) To be catching? (\emph{chuánrǎn}.) A face mask? (\emph{kǒuzhào}.) Mood? (\emph{xīnqíng}.) To look after yourself? (\emph{bǎozhòng}.) A village? (\emph{cūnzi}.) A bridge? (\emph{qiáo}.) A stone? (\emph{shítou}.) A pond? (\emph{chítáng}.) To row a boat? (\emph{huáchuán}.) Something lived through? (\emph{jīnglì}.) Pleasant and cheerful? (\emph{yúkuài}.) A local speciality? (\emph{tèchǎn}.) A souvenir? (\emph{jìniànpǐn}.) To think back on? (\emph{huíyì}.)
\end{tcolorbox}
\section[(dialogue)]{Back From a Trip — once more}

\label{lesson:ZH-R247-review-final-247}



\begin{quote}
Say the Mandarin for something lived through and pleasant and cheerful.
\end{quote}

\subsection*{Your turn}
\begin{itemize}
\raggedright
  \item \emph{Say it:} something lived through, pleasant and cheerful, a local speciality, a souvenir, to think back on
\end{itemize}

\begin{tcolorbox}[breakable,colback=teal!4,colframe=teal!35!black,title={Before you move on}]
Something lived through? (\emph{jīnglì}.) Pleasant and cheerful? (\emph{yúkuài}.) A local speciality? (\emph{tèchǎn}.) A souvenir? (\emph{jìniànpǐn}.) To think back on? (\emph{huíyì}.)
\end{tcolorbox}
